\documentclass[a4paper,14pt]{extarticle}
\usepackage[utf8]{inputenc}
\usepackage[T2A]{fontenc}
\usepackage[english,russian]{babel}
\usepackage{geometry}
\usepackage{float}
\geometry{left=25mm, right=15mm, top=20mm, bottom=20mm}

\begin{document}

\begin{titlepage}
\begin{center}
\vspace*{2cm}
{\Large \textbf{ПОЯСНИТЕЛЬНАЯ ЗАПИСКА}} \\
\vspace{0.5cm}
{\large к техническому проекту} \\
\vspace{0.5cm}
{\huge \textbf{«Motion Visualizer»}} \\
\vspace{0.5cm}
\vfill
г. Ростов-на-Дону, 2026
\end{center}
\end{titlepage}

\section{Обоснование выбора архитектуры}

\subsection{Причины выбора PySide6}
\begin{itemize}
    \item Лицензия LGPL (разрешено коммерческое использование)
    \item Встроенная поддержка OpenGL (QOpenGLWidget)
    \item Кроссплатформенность (Linux/Windows/macOS)
    \item Интеграция с Qt Designer для быстрого прототипирования GUI
\end{itemize}

\subsection{Причины выбора C++ для вычислительного ядра}
\begin{itemize}
    \item Производительность (отсутствие GIL, нативная работа с памятью)
    \item Предсказуемое время выполнения (нет GC-пауз)
    \item Возможность использования SIMD и многопоточности
    \item pybind11 обеспечивает простую интеграцию с Python
\end{itemize}

\subsection{Причины выбора SQLite}
\begin{itemize}
    \item Не требует отдельного сервера
    \item Простая интеграция через SQLAlchemy
    \item Достаточно для объёма данных (логи испытаний)
    \item Переносимость (один файл .db)
\end{itemize}

\section{Описание алгоритмов}

\subsection{Метод Эйлера для интегрирования}
Используется простейший метод численного интегрирования:
\[
x(t + \Delta t) = x(t) + v_x \cdot \Delta t
\]
\[
y(t + \Delta t) = y(t) + v_y \cdot \Delta t
\]
Выбран за линейную сложность и достаточную точность для демонстрации.

\subsection{Алгоритм перетаскивания мышью}
\begin{enumerate}
    \item При нажатии ЛКМ — запоминаем координаты курсора
    \item При движении — вычисляем смещение и добавляем к позиции объекта
    \item При отпускании — прекращаем обновление позиции
\end{enumerate}

\section{Паттерны проектирования}

\begin{table}[H]
\centering
\begin{tabular}{|p{4cm}|p{10cm}|}
\hline
\textbf{Паттерн} & \textbf{Применение} \\
\hline
Command & Кнопки «Старт», «Сброс», «Показать логи» инкапсулированы в методы \\
\hline
Singleton & DatabaseClient (одно соединение с БД) \\
\hline
Bridge & C++ модуль отделён от Python через pybind11 (можно заменить реализацию) \\
\hline
\end{tabular}
\end{table}

\end{document}